\begingroup
\renewenvironment{Proof}{\par}{\par\nobreak\hfill$\blacksquare$\par\vspace{0.1cm}}
\begin{Proof}
We now estimate the principal part. We assert the fact that each function $a_{\sigma\gamma}$ can be extended to be a continuous function on all of $\mathbb{R}^n$. (We already know this to be true for Lipschitz domains since they have the $k$-extension property for any $k$. In fact, by the Tietze extension theorem (consult a topology text), it holds for any domain $\Omega$.)

Now let $\Omega\prime$ be any bounded open domain such that $\Omega$ is compactly contained in $\Omega\prime$. Since each extended $a_{\sigma\gamma}$ ($|\sigma|=|\gamma|=k$) is uniformly continuous on $\Omega\prime$, there exists a nondecreasing \emph{modulus of continuity} function $\omega:[0,\infty)\to[0,\infty)$ satisfying
\begin{equation}\label{eq:garding-common-continuity-modulus-at-zero}
0=\omega(0)=\lim_{\delta\to0+}\omega(\delta)
\end{equation}
and
\begin{equation}\label{eq:garding-common-coefficient-continuity-modulus}
|a_{\sigma\gamma}(\bm{x})-a_{\sigma\gamma}(\bm{y})|\leq\omega(|\bm{x}-\bm{y}|)
\end{equation}
for every $|\sigma|=|\gamma|=k$ and every $\bm{x},\bm{y}\in\Omega\prime$.

Now let $B=B(\bm{x}_0,\delta)$ for some $\bm{x}_0\in\Omega$. We will choose $\delta>0$ later, but for now we assume only that it is sufficiently small so that $B\subset\Omega\prime$. The first step in our estimate of $I_1$ is to do an estimate in the case where $u\in H_0^k(B)$. In this case we have
\begin{equation}\label{eq:garding-frozen-principal-coefficient-splitting}
I_1=I_{11}+I_{12},
\end{equation}
where
\begin{equation}\label{eq:garding-frozen-principal-part}
I_{11}:=\sum_{|\sigma|=|\gamma|=k}\int_{\mathbb{R}^n}a_{\sigma\gamma}(\bm{x}_0)D^\gamma uD^\sigma u\,\mathrm{d}\bm{x},
\end{equation}

\begin{equation}\label{eq:garding-coefficient-variation-part}
I_{12}:=\sum_{|\sigma|=|\gamma|=k}\int_B[a_{\sigma\gamma}(\bm{x})-a_{\sigma\gamma}(\bm{x}_0)]D^\gamma uD^\sigma u\,\mathrm{d}\bm{x}.
\end{equation}
(Note that in the definition of $I_{11}$ we have assumed $u$ is extended by $0$ to all of $\mathbb{R}^n$.)

We can use \eqref{eq:garding-common-coefficient-continuity-modulus} and Hölder\textquotesingle s inequality to get an easy estimate for $I_{12}$:
\begin{equation}\label{eq:garding-coefficient-variation-modulus-bound}
|I_{12}|\leq\omega(\delta)\|u\|_{k,2}^2.
\end{equation}
To estimate $I_{11}$ we use Fourier transforms\[\begin{aligned}
I_{11}&=\sum_{|\sigma|=|\gamma|=k}a_{\sigma\gamma}(\bm{x}_0)\int_{\mathbb{R}^n}D^\gamma u(\bm{x})\overline{D^\sigma u}(\bm{x})\,\mathrm{d}\bm{x}\\
&=\sum_{|\sigma|=|\gamma|=k}a_{\sigma\gamma}(\bm{x}_0)\int_{\mathbb{R}^n}\widehat{D^\gamma u}(\bm{\xi})\overline{\widehat{D^\sigma u}}(\bm{\xi})\,\mathrm{d}\bm{\xi}\\
&=\sum_{|\sigma|=|\gamma|=k}\int_{\mathbb{R}^n}a_{\sigma\gamma}(\bm{x}_0)(i\bm{\xi})^\gamma(-i\bm{\xi})^\sigma|\widehat u|^2\,\mathrm{d}\bm{\xi}\\
&=\sum_{|\sigma|=|\gamma|=k}\int_{\mathbb{R}^n}\bm{\xi}^\sigma a_{\sigma\gamma}(\bm{x}_0)\bm{\xi}^\gamma|\widehat u|^2\,\mathrm{d}\bm{\xi}\\
&\geq\theta\int_{\mathbb{R}^n}|\bm{\xi}|^{2k}|\widehat u|^2\,\mathrm{d}\bm{\xi}.
\end{aligned}\]In the last inequality we have used the uniform ellipticity condition. To continue, we use Theorem~\ref{thm:sobolev-fourier-weighted-homeomorphism} to get\[\begin{aligned}
I_{11}&\geq\theta\int_{\mathbb{R}^n}(1+|\bm{\xi}|^{2k})|\widehat u|^2\,\mathrm{d}\bm{\xi}-\theta\int_{\mathbb{R}^n}|\widehat u|^2\,\mathrm{d}\bm{\xi}\\
&\geq\overline C\|u\|_{k,2}^2-\theta\|u\|_2^2
\end{aligned}\]for some $\overline C>0$ which depends only on $\Omega\prime$.

We now combine the estimates of $I_{11}$ and $I_{12}$ to get an estimate for $I_1$. At this time we assume that $\delta$ is sufficiently small so that $\omega(\delta)\leq\overline C/2$. Then we have
\begin{equation}\label{eq:garding-ball-local-principal-part-bound}
\begin{aligned}
I_1&\geq I_{11}-|I_{12}|\\
&\geq\overline C\|u\|_{k,2}^2-\theta\|u\|_2^2-\omega(\delta)\|u\|_{k,2}^2\\
&\geq\dfrac{\overline C}{2}\|u\|_{k,2}^2-\theta\|u\|_2^2.
\end{aligned}
\end{equation}
We now continue with our estimate of $I_1$ in the case of a general $u\in H_0^k(\Omega)$. The basic idea is to break up $u$ using a partition of unity, so that we can use the previous estimate.

We begin by covering $\overline\Omega$ with a finite collection of balls
\begin{equation}\label{eq:garding-localization-finite-ball-cover}
B_i:=B(\bm{x}_i,\delta_i),\quad i=1,\ldots,M,
\end{equation}
with $\bm{x}_i\in\Omega$ and $\delta_i>0$, selected as in the previous estimate so that $\overline B_i\subset\Omega\prime$. Now let $\psi_i$ be a partition of unity on $\overline\Omega$ subordinate to the covering $\overline B_i$. We then set
\begin{equation}\label{eq:garding-square-normalized-partition}
\phi_i(\bm{x}):=\left(\dfrac{\psi_i^2(\bm{x})}{\sum_{j=1}^M\psi_j^2(\bm{x})}\right)^{1/2}.
\end{equation}
We then have\begin{enumerate}
\item $0\leq\phi_i(\bm{x})\leq1$,
\item $\phi_i\in C^\infty(B_i\cap\overline\Omega)$,
\item $\sum_{i=1}^M\phi_i^2(\bm{x})=1$ for each $\bm{x}\in\Omega$, and
\item $u_i:=u\phi_i\in H_0^k(B_i)$.
\end{enumerate}This can be used to write\[\begin{aligned}
I_1&=\sum_{|\sigma|=|\gamma|=k}\int_\Omega a_{\sigma\gamma}(\bm{x})D^\sigma uD^\gamma u\,\mathrm{d}\bm{x}\\
&=\sum_{i=1}^M\sum_{|\sigma|=|\gamma|=k}\int_\Omega a_{\sigma\gamma}(\bm{x})\phi_iD^\sigma u\phi_iD^\gamma u\,\mathrm{d}\bm{x}\\
&=\sum_{i=1}^M\sum_{|\sigma|=|\gamma|=k}\int_\Omega a_{\sigma\gamma}(\bm{x})D^\sigma(\phi_i u)D^\gamma(\phi_i u)\,\mathrm{d}\bm{x}\\
&\quad+\sum_{i=1}^M\sum_{|\sigma|=|\gamma|=k}\int_\Omega a_{\sigma\gamma}(\bm{x})[\phi_iD^\sigma u-D^\sigma(\phi_i u)]\phi_iD^\gamma u\,\mathrm{d}\bm{x}\\
&\quad+\sum_{i=1}^M\sum_{|\sigma|=|\gamma|=k}\int_\Omega a_{\sigma\gamma}(\bm{x})[\phi_iD^\gamma u-D^\gamma(\phi_i u)]D^\sigma(\phi_i u)\,\mathrm{d}\bm{x}\\
&\geq\sum_{i=1}^M\sum_{|\sigma|=|\gamma|=k}\int_\Omega a_{\sigma\gamma}(\bm{x})D^\sigma u_iD^\gamma u_i\,\mathrm{d}\bm{x}\\
&\quad-C\|u\|_{k,2}\|u\|_{k-1,2}.
\end{aligned}\]We can now use the previous estimate for each $u_i\in H_0^k(B_i)$ to get\[\begin{aligned}
I_1&\geq\dfrac{\overline C}{2}\sum_{i=1}^M\|u_i\|_{k,2}^2-C\|u\|_2^2-C\|u\|_{k,2}\|u\|_{k-1,2}\\
&=\dfrac{\overline C}{2}\sum_{i=1}^M\sum_{|\alpha|\leq k}\int_\Omega|D^\alpha(\phi_i u)|^2\,\mathrm{d}\bm{x}\\
&\quad-C[\|u\|_2^2+\|u\|_{k,2}\|u\|_{k-1,2}]\\
&\geq\dfrac{\overline C}{2}\sum_{|\alpha|\leq k}\int_\Omega\sum_{i=1}^M(\phi_i^2)|D^\alpha u|^2\,\mathrm{d}\bm{x}\\
&\quad-C[\|u\|_2^2+\|u\|_{k,2}\|u\|_{k-1,2}]\\
&=\dfrac{\overline C}{2}\|u\|_{k,2}^2-C[\|u\|_2^2+\|u\|_{k,2}\|u\|_{k-1,2}]
\end{aligned}\]Thus, using \eqref{eq:garding-elementary-young-product-inequality} and \eqref{eq:ehrling-compact-intermediate-norm-bound} we get
\begin{equation}\label{eq:garding-global-principal-part-lower-bound}
I_1\geq\dfrac{\overline C}{4}\|u\|_{k,2}^2-C\|u\|_2^2.
\end{equation}
Finally we combine this estimate with \eqref{eq:garding-final-lower-order-estimate} to get\[\begin{aligned}
B[u,u]&=I_1+I_2\\
&\geq\left(\dfrac{\overline C}{4}-\varepsilon\right)\|u\|_{k,2}^2-C(\varepsilon)\|u\|_2^2\\
&:=c_3\|u\|_{k,2}^2-\lambda_G\|u\|_2^2,
\end{aligned}\]where in defining $c_3$ and $\lambda_G$ we have taken $\varepsilon$ to be sufficiently small, say $\varepsilon=\overline C/8$.
\end{Proof}
\endgroup
